\chapter{Genel Sınava Hazırlık}
\label{ch:review}

\begin{hedefler}
Bu bölümün sonunda:
\begin{itemize}
  \item araştırma sorusu, değişken türü ve gözlem yapısından uygun yönteme ulaşabilecek,
  \item örnekleme, betimleme, tahmin ve hipotez testi kavramlarını aynı problemde birleştirebilecek,
\ifimoders
  \item \tstat testleri, ki-kare, korelasyon ve basit regresyon çıktılarını varsayımlarıyla yorumlayabilecek,
\else
  \item \tstat testleri, ANOVA, ki-kare, korelasyon ve regresyon çıktılarını varsayımlarıyla yorumlayabilecek,
\fi
  \item hesabı, belirsizliği, etki büyüklüğünü ve bağlamsal sonucu içeren bütünlüklü sınav yanıtı yazabileceksiniz.
\end{itemize}
\end{hedefler}

\begin{prerequisitebox}
\ifimoders
Bölüm~\ref{ch:research}--\ref{ch:correlation-regression}'teki temel
kavramların ilk kez öğrenilmiş olması beklenir.
\else
Bölüm 1--19'daki temel kavramların ilk kez öğrenilmiş olması beklenir.
\fi
Bu
bölüm yeni bir yöntem öğretmekten çok kavramlar arasındaki bağlantıları
pekiştirir. Veri üretiminden yorumlamaya uzanan bu bütüncül yaklaşım GAISE
önerileriyle uyumludur \citep{gaise2016}.
\end{prerequisitebox}

\begin{hazirlik}
Bir analiz sorusunu gördüğünüzde yöntem adını düşünmeden önce yazmanız gereken
dört bilgiyi sıralayın: hedef, değişken türü, gözlem yapısı ve belirsizlik.
\end{hazirlik}

\begin{definitionbox}
\textbf{Bütünlüklü istatistiksel yanıt}; hedef parametreyi, yöntem seçiminin
gerekçesini, varsayım ve tasarım sınırlarını, sayısal sonucu, belirsizliği ve
bağlamsal yorumu birlikte içerir.
\end{definitionbox}

\ifimoders
Bu bölüm kitabın önceki yöntemlerini birlikte kullanır. Tahmin ve aralıklar
için Bölüm~\ref{ch:point-estimation}--\ref{ch:confidence-intervals}; testler
için Bölüm~\ref{ch:hypothesis}--\ref{ch:categorical}; korelasyon ve basit
doğrusal regresyon için Bölüm~\ref{ch:correlation-regression} hızlı başvuru
noktalarıdır.
\else
Bu bölüm kitabın önceki yöntemlerini birlikte kullanır. Tahmin ve aralıklar
için Bölüm~\ref{ch:point-estimation}--\ref{ch:confidence-intervals}; testler
için Bölüm~\ref{ch:hypothesis}--\ref{ch:categorical}; ilişki modelleri için
Bölüm~\ref{ch:correlation-regression}--\ref{ch:multiple-regression} ilişki
modelleri; Bölüm~\ref{ch:anova} ise üç veya daha fazla bağımsız ortalamanın
karşılaştırılması için hızlı başvuru noktalarıdır.
\fi

\section{Kavramsal kontrol listesi}

\begin{enumerate}
  \item Evren, örneklem, parametre ve istatistiği ayırt edin.
  \item Örnekleme yöntemini ve olası yanlılık kaynaklarını belirleyin.
  \item Değişken türüne uygun grafik ve özetleri seçin.
  \item Standart sapma ile standart hatayı ayırt edin.
  \item Örnekleme dağılımı ve MLT arasındaki ilişkiyi açıklayın.
  \item Nokta tahmini ile güven aralığının farklı rollerini belirtin.
  \item \pval-değerini, hata türlerini ve test gücünü doğru yorumlayın.
  \item Araştırma tasarımına uygun \tstat testini seçin.
\ifimoders\else
  \item Üç veya daha fazla bağımsız ortalama için ANOVA'yı ve çoklu karşılaştırmaları değerlendirin.
\fi
  \item Kategorik değişkenler için çapraz tablo ve ki-kare testini değerlendirin.
  \item Korelasyon, regresyon ve nedensellik kavramlarını ayırın.
\end{enumerate}

\section{Yöntem seçim özeti}
\label{sec:method-summary}

\begin{table}[H]
\centering
\small
\begin{tabularx}{\textwidth}{@{}>{\raggedright\arraybackslash}X>{\raggedright\arraybackslash}X@{}}
\toprule
Araştırma hedefi & Temel yöntem \\
\midrule
Tek ortalamayı sabit değerle karşılaştırma & Tek örneklem \tstat testi \\
Tek ortalama veya oranı tahmin etme & Standart hata ve güven aralığı \\
İki bağımsız grup ortalamasını karşılaştırma & Bağımsız örneklem Welch \tstat testi \\
Aynı bireylerin iki ölçümünü karşılaştırma & Eşleştirilmiş örneklem \tstat testi \\
\ifimoders\else
Üç veya daha fazla bağımsız grup ortalamasını karşılaştırma & Tek yönlü ANOVA veya Welch ANOVA \\
\fi
Tek kategorik dağılımı kuramsal oranlarla karşılaştırma & Ki-kare uyum iyiliği testi \\
İki kategorik değişkenin ilişkisini inceleme & Çapraz tablo ve ki-kare bağımsızlık testi \\
İki nicel değişkenin doğrusal ilişkisini ölçme & Pearson korelasyonu \\
Nicel yanıtı tek açıklayıcıyla modelleme & Basit doğrusal regresyon \\
\bottomrule
\end{tabularx}
\end{table}

\section{Yöntem seçme akışı}

Bir sınav sorusunda yöntemi doğrudan formüllerden aramak yerine önce yanıt
değişkeninin türü ve gözlemlerin ilişkisi belirlenmelidir.

\begin{figure}[H]
\centering
\resizebox{\textwidth}{!}{%
\begin{tikzpicture}[every node/.style={font=\scriptsize,align=center}]
  \node[draw,rounded corners,fill=PrismLight,minimum width=3.4cm,
    minimum height=0.8cm] (start) at (0,0) {Yanıt değişkeni nedir?};
  \node[draw,rounded corners,fill=PrismNoteBlue,minimum width=3.1cm,
    minimum height=0.8cm] (num) at (-3.5,-1.7) {Nicel};
  \node[draw,rounded corners,fill=PrismResearchGreen,minimum width=3.1cm,
    minimum height=0.8cm] (cat) at (3.5,-1.7) {Kategorik};
  \node[draw,fill=white,minimum width=3.3cm,minimum height=1.15cm]
    (one) at (-5.2,-3.7) {Tek ortalama\\tahmin veya tek örneklem t};
  \node[draw,fill=white,minimum width=3.3cm,minimum height=1.15cm]
    (two) at (-1.7,-3.7) {\ifimoders İki ortalama\\bağımsız/eşleştirilmiş t\else Grup ortalamaları\\2 grup: t; 3+: ANOVA\fi};
  \node[draw,fill=white,minimum width=3.3cm,minimum height=1.15cm]
    (assoc) at (1.7,-3.7) {İki nicel değişken\\korelasyon/regresyon};
  \node[draw,fill=white,minimum width=3.3cm,minimum height=1.15cm]
    (chi) at (5.2,-3.7) {Frekans dağılımı\\uyum iyiliği/ki-kare};
  \draw[->,PrismBlue,thick] (start)--(num);
  \draw[->,PrismBlue,thick] (start)--(cat);
  \draw[->,PrismBlue] (num)--(one);
  \draw[->,PrismBlue] (num)--(two);
  \draw[->,PrismBlue] (num)--(assoc);
  \draw[->,PrismBlue] (cat)--(chi);
\end{tikzpicture}
}
\caption{Araştırma sorusundan temel yönteme geçiş için kısa karar akışı.}
\label{fig:method-decision-flow}
\end{figure}

Şekil~\ref{fig:method-decision-flow}'deki akış son karar değildir. Bağımsızlık,
örnekleme yöntemi, dağılım biçimi,
beklenen frekanslar ve ölçümün niteliği ayrıca kontrol edilir. Yöntem seçimi
bir değişkenin yazılımda sayısal görünmesine değil, araştırmadaki anlamına
dayanır.

\section{Çözümlü örnek: yöntem seçimi ve yorum}

Bir araştırmacı aynı 24 öğrencinin eğitim öncesi ve sonrası kaygı puanlarını
karşılaştırmak istemektedir. Son--ön farkının ortalaması $-3{,}2$, farkların
standart sapması 6,0 bulunmuştur.

\begin{enumerate}
  \item Ölçümler aynı öğrencilere ait olduğu için eşleştirilmiş örneklem \tstat testi seçilir.
  \item Hedef parametre son--ön farklarının evren ortalaması $\mu_D$'dir.
  \item Standart hata $6/\sqrt{24}=1{,}225$ ve test değeri
  $t=-3{,}2/1{,}225=-2{,}61$'dir.
  \item $sd=23$ için çift yönlü \pval-değeri yaklaşık 0,016'dır; $H_0:\mu_D=0$ reddedilir.
  \item Yüzde 95 güven aralığı yaklaşık $[-5{,}73;-0{,}67]$, etki büyüklüğü
  $d_z=-3{,}2/6=-0{,}53$'tür.
\end{enumerate}

Negatif işaret son ölçüm kaygı puanının ortalama olarak daha düşük olduğunu
gösterir. Kontrol grubu olmayan bir ön test--son test tasarımında azalmanın
yalnız eğitime bağlı olduğu söylenemez.

\section{Çözümlü karma problemler}

\subsection*{Problem 1: Örnekleme, değişken türü ve genelleme}

Bir üniversite, öğrencilerin haftalık çalışma süresi ile sınav başarısı
arasındaki ilişkiyi incelemek için kütüphaneden çıkan ilk 200 gönüllüye anket
uygulamıştır. Çalışma süresi saat, başarı ise başarılı/başarısız olarak
kaydedilmiştir. Evreni, örneklemi, değişken türlerini, temel yanlılık riskini
ve uygun ilk analizi belirleyin.

\textbf{Çözüm.} Hedef evren üniversitedeki araştırma kapsamına giren bütün
öğrenciler, örneklem ankete katılan 200 gönüllüdür. Çalışma süresi sürekli
nicel, başarı durumu ikili kategorik değişkendir. Yalnız kütüphaneden çıkan
gönüllülerin seçilmesi kapsam ve gönüllülük yanlılığı doğurabilir; düzenli
kütüphane kullanan öğrenciler evrenin geri kalanından farklı olabilir.

İlk aşamada başarı gruplarındaki çalışma süresi dağılımları yan yana kutu veya
nokta grafikleriyle karşılaştırılabilir. İki bağımsız başarı grubunun ortalama
çalışma süreleri incelenecekse koşullar altında Welch \tstat testi düşünülebilir.
Ancak gözlemsel ve yanlı seçilmiş veri, çalışma süresinin başarıya neden
olduğu sonucunu desteklemez.

\subsection*{Problem 2: Örnekleme dağılımı ve olasılık}

Ortalaması 60, standart sapması 12 olan bir evrenden bağımsız $n=36$
gözlem seçilmektedir. Örneklem ortalamasının beklenen değerini ve standart
hatasını bulun. Merkezi limit yaklaşımıyla $\P(\bar X>64)$ olasılığını
hesaplayın.

\textbf{Çözüm.}
\[
\E[\bar X]=60,
\qquad
\SE(\bar X)=\frac{12}{\sqrt{36}}=2.
\]
64 değeri örnekleme dağılımı merkezinin
\[
z=\frac{64-60}{2}=2
\]
standart hata üzerindedir. Standart normal dağılımdan
$\P(Z>2)\approx0{,}0228$ bulunur. Yani aynı örnekleme süreci tekrarlandığında
örneklem ortalamasının 64'ü aşması yaklaşık yüzde 2,28 olasılıklıdır. Bu
hesap tek tek öğrenci puanlarının 64'ü aşma olasılığı değildir.

\subsection*{Problem 3: Güven aralığı ve örneklem planlama}

Bağımsız 49 öğrencinin ortalama puanı 72, standart sapması 14'tür.
$t_{0{,}975,48}=2{,}01$ kullanarak yüzde 95 güven aralığını bulun. Standart
sapma yaklaşık 14 kabul edilirse yüzde 95 güven düzeyinde hata payını 2 puana
indirmek için normal yaklaşımla gereken örneklem hacmini hesaplayın.

\textbf{Çözüm.} Standart hata $14/\sqrt{49}=2$ ve hata payı
$2{,}01(2)=4{,}02$'dir. Güven aralığı
\[
72\pm4{,}02=[67{,}98;76{,}02]
\]
olur. Hedef hata payı için
\[
n\approx\left(\frac{1{,}96(14)}{2}\right)^2
=188{,}24
\]
bulunur ve örneklem hacmi yukarı yuvarlanarak en az 189 seçilir. Bu plan basit
rastgele örnekleme yaklaşımına dayanır; beklenen yanıtlamama veya kümeleşme
varsa başlangıç sayısı artırılmalıdır.

\subsection*{Problem 4: Tek örneklem testi, aralık ve etki büyüklüğü}

Bir ölçeğin referans ortalaması 50'dir. $n=25$, $\bar x=54$ ve $s=10$
değerleri için çift yönlü tek örneklem \tstat testini, yüzde 95 güven aralığını ve
Cohen'in $d$ değerini hesaplayın.

\textbf{Çözüm.} Standart hata $10/\sqrt{25}=2$ ve
\[
t=\frac{54-50}{2}=2{,}00,\qquad sd=24
\]
olur. Çift yönlü \pval-değeri yaklaşık 0,057'dir; yüzde 5 düzeyinde $H_0$
reddedilemez. $t_{0{,}975,24}=2{,}064$ ile ortalama güven aralığı
\[
54\pm2{,}064(2)=[49{,}87;58{,}13]
\]
olup referans değer 50'yi içerir. Etki büyüklüğü
$d=(54-50)/10=0{,}40$'tır.

Sonuç ``fark yoktur'' biçiminde yorumlanmaz. Örneklemde dört puanlık fark
gözlenmiş, ancak aralık küçük negatif farklara yakın değerlerden sekiz puanın
üzerindeki pozitif farklara kadar uzanmıştır. Çalışmanın hassasiyeti sınırlıdır.

\subsection*{Problem 5: Bağımsız iki grup için Welch testi}

A grubunda $n_1=40$, $\bar x_1=78$, $s_1=10$; B grubunda
$n_2=35$, $\bar x_2=72$, $s_2=14$ bulunmuştur. Yazılım Welch testi için
$sd=60{,}60$ vermektedir. Test istatistiğini, yüzde 95 güven aralığını ve
yaklaşık etki büyüklüğünü hesaplayıp yorumlayın.

\textbf{Çözüm.} Ortalama farkı 6 ve standart hata
\[
SE=\sqrt{\frac{10^2}{40}+\frac{14^2}{35}}
=\sqrt{8{,}10}\approx2{,}85
\]
olduğundan
\[
t=\frac{6}{2{,}85}\approx2{,}11.
\]
Çift yönlü \pval-değeri yaklaşık 0,039'dur. Uygun kritik değerle yüzde 95 fark
aralığı yaklaşık $[0{,}31;11{,}69]$ bulunur. Birleşik standart sapma yaklaşık
12,03 olduğundan betimsel standartlaştırılmış fark $d\approx0{,}50$'dir.

A grubunun ortalamasının daha yüksek olduğuna ilişkin kanıt vardır; ancak
makul farklar yaklaşık 0,3 ile 11,7 puan arasında geniştir. Gruplara rastgele
atama yapılmadıysa fark yöntemlerin nedensel etkisi olarak yorumlanamaz.

\subsection*{Problem 6: Ki-kare ve hücre örüntüsü}

İki programdaki başarı durumu aşağıdaki gibidir:
\[
\begin{array}{c|rr|r}
 & \text{Başarılı} & \text{Başarısız} & \text{Toplam}\\ \hline
A&36&24&60\\
B&24&36&60\\ \hline
\text{Toplam}&60&60&120
\end{array}
\]
Pearson ki-kare değerini, serbestlik derecesini ve Cramér $V$ değerini
hesaplayarak sonucu yorumlayın.

\textbf{Çözüm.} Bütün hücrelerin beklenen frekansı $60(60)/120=30$'dur.
Her hücrenin katkısı $(6^2)/30=1{,}20$ olduğundan
\[
\chi^2=4(1{,}20)=4{,}80,
\qquad sd=1,
\qquad p\approx0{,}028.
\]
Cramér katsayısı
\[
V=\sqrt{\frac{4{,}80}{120}}=0{,}20
\]
olur. A programında başarı oranı yüzde 60, B programında yüzde 40'tır.
Program ile başarı durumu arasında ilişkiye yönelik kanıt vardır. İlişkinin
yönü yüzdelerden, büyüklüğü $V$ değerinden okunur; nedensellik için atama ve
olası karıştırıcılar ayrıca değerlendirilir.

\subsection*{Problem 7: Korelasyon ve regresyonu birleştirme}

Bir çalışmada $n=50$, $r=0{,}50$, $\bar x=4$, $s_X=2$,
$\bar y=60$ ve $s_Y=8$ bulunmuştur. Regresyon doğrusunu, $R^2$ değerini,
korelasyon testini ve $X=7$ için öngörüyü hesaplayın.

\textbf{Çözüm.} Eğim ve sabit terim
\[
b_1=0{,}50\frac{8}{2}=2,
\qquad b_0=60-2(4)=52
\]
olduğundan regresyon doğrusu $\hat Y=52+2X$'tir. Açıklanan varyans oranı
$R^2=0{,}25$ ve $X=7$ için öngörü $\hat y=66$'dır. Korelasyon testi
\[
t=0{,}50\sqrt{\frac{48}{1-0{,}25}}=4{,}00,
\qquad sd=48
\]
sonucunu verir; çift yönlü \pval-değeri 0,001'den küçüktür.

Pozitif doğrusal ilişkiye ilişkin güçlü istatistiksel kanıt vardır ve model
örneklemdeki $Y$ değişkenliğinin yüzde 25'ini açıklamaktadır. Eğim, $X$'teki
bir birimlik farkla ilişkili ortalama iki birimlik $Y$ farkıdır; gözlemsel
çalışmada nedensel artış olarak sunulamaz.

\section{Yazılım uygulaması: hesapları denetleme}

\Needspace{22\baselineskip}
\begin{softwarebox}
\begin{lstlisting}
from scipy import stats
import numpy as np

n, mean_diff, sd_diff = 24, -3.2, 6.0
se = sd_diff / np.sqrt(n)
t_value = mean_diff / se
p_value = 2 * stats.t.sf(abs(t_value), df=n-1)
critical = stats.t.ppf(.975, df=n-1)
ci = (mean_diff - critical*se,
      mean_diff + critical*se)

print(t_value, p_value, ci)
\end{lstlisting}
Yazılım doğrulaması yöntemin gerekçesinin yerini tutmaz. İşaretin yönü,
farkın ``son--ön'' ya da ``ön--son'' tanımına göre yorumlanmalıdır.
\end{softwarebox}

\begin{mistakebox}
Sınav çözümünde yalnız ``$p<0{,}05$, anlamlıdır'' yazmak yeterli değildir.
Hedef parametre, test türü, farkın yönü, belirsizlik ve bağlamsal sonuç açıkça
belirtilmelidir.
\end{mistakebox}

\begin{assumptionbox}
Bir yöntemin formülünü uygulamadan önce gözlemlerin bağımsızlığı, değişken
türü, grup veya eşleşme yapısı ve dağılım koşulları kontrol edilmelidir.
Varsayım değerlendirmesi sonuç görüldükten sonra yapılan bir formalite değildir.
\end{assumptionbox}

\section{Sınav yanıtı yazma ilkesi}

Her çözümde hedef parametreyi, kullanılan yöntemi, gerekli koşulları,
hesaplanan istatistiği ve bağlamsal yorumu açıkça belirtin. Yalnız formül veya
yazılım çıktısı vermek tam bir istatistiksel yanıt değildir.

Tam puana yaklaşan kısa bir yanıt şu sırayı izler:

\begin{enumerate}
  \item \textbf{Tanımla:} Hedef evreni, parametreyi ve değişkenleri yazın.
  \item \textbf{Seç:} Yöntemi veri yapısına dayanarak gerekçelendirin.
  \item \textbf{Kontrol et:} Bağımsızlık ve yönteme özgü koşulları belirtin.
  \item \textbf{Hesapla:} Tahmin, standart hata, test istatistiği veya aralığı gösterin.
  \item \textbf{Karar ver:} \pval-değerini önceden belirlenen $\alpha$ ile karşılaştırın.
  \item \textbf{Yorumla:} Yönü, büyüklüğü, belirsizliği ve tasarım sınırını bağlamda açıklayın.
\end{enumerate}

\begin{prismbox}[PrismResearchGreen]{Son kontrol}
Bir yanıtı teslim etmeden önce dört soruyu sorun: Hedef büyüklüğü belirttim
mi? Yöntemi veri yapısına göre seçtim mi? Belirsizliği gösterdim mi? Sonucu
araştırma bağlamında ve yöntemin sınırlarıyla birlikte yorumladım mı?
\end{prismbox}

\section{Literatür verisiyle yazılım uygulamaları}
\label{sec:spss-b14}

\textbf{Amaç ve veri.} \texttt{b14.csv}, aynı literatür verisiyle
betimleme, aralık tahmini, test ve ilişki analizlerini bir raporda
birleştirir \citep{cortez2008data}. Bu bir yöntem seçimi ve raporlama
alıştırmasıdır; çok sayıda testten yalnız anlamlı olan seçilmez.
Üç ayrı soru vardır: G3 ortalaması 10'dan farklı mı; kaynak cinsiyet
kategorisi ile 10 puan eşiğini karşılama durumu ilişkili mi;
G1 ve G3 arasında doğrusal ilişki var mı? Sıfır hipotezleri sırasıyla
$\mu_{G3}=10$, kategorik değişkenlerin bağımsızlığı ve $\rho=0$'dır.
Tek örneklem ve korelasyon testleri iki yönlüdür; kategorik analiz
süreklilik düzeltmesiz Pearson ki-kare testidir. Her test için
anlamlılık düzeyi 0,05'tir; çoklu test düzeltmesi yapılmamıştır.

\textbf{Ortak veri.} Karşılaştırmanın dosyası
\nolinkurl{companion/spss/csv/b14.csv} olup 395 satır ve 10 sütun
içerir. Ana analizler tüm kayıtları kullanır; sıfır yıl sonu notu
bulunan 38 kayıt eksik sayılmaz. \texttt{sex} için 1=F ve 2=M;
\texttt{pass10} için 0: G3$<10$, 1: G3$\geq10$ tanımlıdır.
B06'nın seçim göstergeleri ve ağırlığı kullanılmaz. Bu karşılaştırma
Excel içe aktarmanın doğrulandığı anlamına gelmez. Kodları
\texttt{companion} klasörünü içeren paket kökünden çalıştırın;
veri sözlüğü için bkz. sayfa~\pageref{sec:spss-guide}.

\subsection{IBM SPSS uygulaması}
\textbf{Başlamadan önce.} Aşağıdaki veri açma bloğunu
\texttt{EXECUTE} satırı dahil çalıştırın; sonra menü veya analiz
komutlarıyla devam edin. Filtre, ağırlık ve dosya bölme kapalı
olmalıdır. Dosya açılamazsa \texttt{/FILE} yolunu düzeltip yeniden
açın; önceki bölümün etkin verisiyle devam etmeyin.

\begin{spssadimlar}
\spssadim{Önce temel araştırma sorusunu yazın. \emph{Analyze $\to$ Descriptive Statistics $\to$ Descriptives} yolunda \texttt{g1} ve \texttt{g3} için temel özetleri alın.}
\spssadim{\emph{Analyze $\to$ Descriptive Statistics $\to$ Explore} yolunda \texttt{g3} için kutu grafiği ve yüzde 95 ortalama güven aralığını alın.}
\spssadim{\emph{Analyze $\to$ Compare Means $\to$ One-Sample T Test} yolunda \texttt{g3} ve 10 referans değeriyle testi çalıştırın. Yüzde 95 güven düzeyini seçin; \emph{Two-Sided p} sütununu okuyun. Menü bazı sürümlerde \emph{Compare Means and Proportions} adını taşır.}
\spssadim{\emph{Analyze $\to$ Descriptive Statistics $\to$ Crosstabs} yolunda \texttt{sex} satır, \texttt{pass10} sütun olsun; ki-kare, Cramér $V$, gözlenen/beklenen sayımlar ve satır yüzdelerini isteyin. Karşılaştırmada \emph{Pearson Chi-Square} satırını kullanın.}
\spssadim{\emph{Analyze $\to$ Correlate $\to$ Bivariate} yolunda \texttt{g1} ve \texttt{g3} için \emph{Pearson}, \emph{Two-tailed} seçip çalıştırın.}
\spssadim{Yalnız duyarlılık incelemesi için \emph{Data $\to$ Select Cases} yolunda \texttt{g3>0} koşuluyla kayıtları filtreleyin; silmeyin. \emph{Descriptives} ile \texttt{g3} özetini alın; ardından \emph{All cases} ile filtreyi kapatıp aynı özeti yeniden alın. Sırayla 357 ve 395 kayıt görünmelidir.}
\end{spssadimlar}

{\raggedright
\noindent\textbf{Komutların görevi.} Ana analizler farklı soruların
çıktılarıdır. \texttt{TEMPORARY} ve \texttt{SELECT IF}, yalnız ilk
izleyen betimsel işlemde pozitif notları kullanır. İkinci
\texttt{DESCRIPTIVES}, tüm kayıtlara dönüşü kontrol eder.\par}

\par\noindent\begin{minipage}{\linewidth}
\textbf{Uygulama kodu:} \texttt{b14}. Veri açma ve tanımlama:

\begin{lstlisting}[style=prismSPSS]
SET DECIMAL=DOT.
GET DATA /TYPE=TXT
 /FILE='companion/spss/csv/b14.csv'
 /ENCODING='UTF8'
 /ARRANGEMENT=DELIMITED /DELCASE=LINE
 /DELIMITERS="," /QUALIFIER='"'
 /FIRSTCASE=2
 /VARIABLES=id F8.0 school F8.0 sex F8.0 age F8.0
 studytime F8.0 absences F8.0 g1 F8.0 g2 F8.0
 g3 F8.0 pass10 F8.0.
FILTER OFF.
WEIGHT OFF.
SPLIT FILE OFF.
VARIABLE LABELS
 g1 'Ilk donem matematik notu'
 g3 'Yil sonu matematik notu'
 sex 'Kaynak verideki cinsiyet kategorisi'
 pass10 'Yil sonu notu en az 10'.
VALUE LABELS
 sex 1 'F' 2 'M'
 /pass10 0 '10 puanin alti' 1 '10 ve uzeri'.
VARIABLE LEVEL
 id school sex pass10 (NOMINAL)
 studytime (ORDINAL)
 age absences g1 g2 g3 (SCALE).
FORMATS g1 g3 (F14.9).
EXECUTE.
\end{lstlisting}
\end{minipage}\par

\par\noindent\begin{minipage}{\linewidth}
\textbf{Ana analizler.} Veri açıldıktan sonra çalıştırın.

\begin{lstlisting}[style=prismSPSS]
DESCRIPTIVES VARIABLES=g1 g3
 /STATISTICS=MEAN STDDEV MIN MAX.
EXAMINE VARIABLES=g3 /PLOT=BOXPLOT
 /STATISTICS=DESCRIPTIVES /CINTERVAL=95.
T-TEST /TESTVAL=10 /VARIABLES=g3
 /CRITERIA=CI(.95).
CROSSTABS /TABLES=sex BY pass10
 /FORMAT=AVALUE TABLES
 /STATISTICS=CHISQ PHI
 /CELLS=COUNT EXPECTED ROW.
CORRELATIONS /VARIABLES=g1 g3
 /PRINT=TWOTAIL /MISSING=PAIRWISE.
\end{lstlisting}
\end{minipage}\par

\par\noindent\begin{minipage}{\linewidth}
\textbf{Ayrı duyarlılık özeti ve geri dönüş.} Bu bloğu bütün olarak
çalıştırın. \texttt{TEMPORARY} ile ilk \texttt{DESCRIPTIVES}
arasına \texttt{EXECUTE} veya başka bir analiz komutu eklemeyin.

\begin{lstlisting}[style=prismSPSS]
TEMPORARY.
SELECT IF (g3>0).
DESCRIPTIVES VARIABLES=g3
 /STATISTICS=MEAN STDDEV MIN MAX.
DESCRIPTIVES VARIABLES=g3
 /STATISTICS=MEAN STDDEV MIN MAX.
\end{lstlisting}
\end{minipage}\par

\begin{outputguidebox}
\textbf{Paylaşılan SPSS çıktısıyla doğrulanan değerler.}
G3 için 395 geçerli ve sıfır eksik kayıt vardır. Ortalama
10,415189873 ve standart sapma 4,5814426110'dur. Explore'da
ortalama için yüzde 95 aralık yaklaşık $[9{,}961992;\ 10{,}868388]$
görülür. Tek örneklem testinde $t(394)=1{,}801$ ve iki yönlü
$p=0{,}072$ vardır. Renkli yardım alanında tek yönlü 0,036'nın
anlamlı gösterilmesi bu uygulamanın kararını değiştirmez;
\emph{Two-Sided p} sütunu esas alınır. Test tablosundaki aralık,
ortalamanın değil 10'dan farkın aralığıdır.

Çapraz tabloda F için $(75,133)$, M için $(55,132)$ sayımları;
Pearson satırında $\chi^2(1)=1{,}970$, $p=0{,}160$ ve
\emph{Symmetric Measures} tablosunda Cramér $V=0{,}071$ vardır.
En küçük beklenen sayım 61,54'tür; 5'in altında beklenen hücre
yoktur. Süreklilik düzeltmesindeki 0,195 ve Fisher testindeki
iki yönlü 0,165 farklı yöntemlerin $p$ değerleridir.
Korelasyon tablosu $r=0{,}801$, $p<0{,}001$ ve $n=395$ verir.

Son iki betimsel tablo sırayla 357 ve 395 kayıt gösterir.
İlkinde ortalama 11,523809524, standart sapma 3,2277969402;
ikincisinde yeniden ortalama 10,415189873 bulunmuştur.
Böylece geçici seçimden tüm kayıtlara dönüş de kontrol edilmiştir.
Ortalamadaki artış sıfırların silinmesini haklı çıkarmaz;
pozitif notlu alt grup, bütün kayıtlardan farklı bir hedeftir.
\end{outputguidebox}

\par\noindent\begin{minipage}{\linewidth}
\subsection{R uygulaması}
\label{sec:r-gercek-b14}

Blokları sırayla ve çok satırlı komutları bütünüyle çalıştırın.
\texttt{pozitif}, yalnız duyarlılık incelemesinde kullanılacak ayrı
alt gruptur; ana analizler \texttt{veri} üzerinden yapılır.

\begin{lstlisting}[style=prismR]
veri <- read.csv("companion/spss/csv/b14.csv")
stopifnot(nrow(veri) == 395, ncol(veri) == 10,
  !anyNA(veri), !anyDuplicated(veri$id),
  all(veri$sex %in% c(1, 2)),
  all(veri$pass10 %in% c(0, 1)),
  all(veri$pass10 == as.integer(veri$g3 >= 10)))
betimsel <- function(degerler) {
  c(n=length(degerler), ortalama=mean(degerler),
    ss=sd(degerler), minimum=min(degerler),
    maksimum=max(degerler))
}
print(c(satir=nrow(veri), sutun=ncol(veri),
  eksik=sum(is.na(veri)), sifir_not=sum(veri$g3 == 0)))
print(rbind(G1=betimsel(veri$g1), G3=betimsel(veri$g3)),
      digits=15)
\end{lstlisting}
\end{minipage}\par

\textbf{Dosya yolu sorunu yaşarsanız.} İlk satır yerine aşağıdaki
komutla \texttt{b14.csv} dosyasını seçin; sonra kontrol bloğundan
devam edin.
\begin{lstlisting}[style=prismR]
veri <- read.csv(file.choose())
\end{lstlisting}

\par\noindent\begin{minipage}{\linewidth}
\textbf{Ortalama, fark ve test.} İki aralığın etiketlerini ayırt edin.

\begin{lstlisting}[style=prismR]
test <- t.test(veri$g3, mu=10,
  alternative="two.sided", conf.level=.95)
print(c(t=unname(test$statistic),
  df=unname(test$parameter), p_iki_yonlu=test$p.value,
  ortalama_fark=mean(veri$g3)-10,
  ortalama_alt=test$conf.int[1],
  ortalama_ust=test$conf.int[2],
  fark_alt=test$conf.int[1]-10,
  fark_ust=test$conf.int[2]-10), digits=15)
\end{lstlisting}

\textbf{Çıktıyı okuma.} R ekranında iki yönlü
$p=0{,}0724481981327951$ bulunmuştur. Ortalama aralığının
her iki sınırından 10 çıkarılınca SPSS'in fark aralığı elde edilir.
\end{minipage}\par

\par\noindent\begin{minipage}{\linewidth}
\textbf{Kategorik ilişki.} \texttt{correct=FALSE}, süreklilik
düzeltmesini kapatır; satır ve sütun sıraları açıkça sabitlenir.

\begin{lstlisting}[style=prismR]
tablo <- table(
  sex=factor(veri$sex, levels=c(1, 2),
             labels=c("F", "M")),
  pass10=factor(veri$pass10, levels=c(0, 1),
                labels=c("10_alti", "10_ve_uzeri")))
ki_kare <- chisq.test(tablo, correct=FALSE)
print(addmargins(tablo))
print(100*prop.table(tablo, margin=1), digits=15)
print(ki_kare$expected, digits=15)
print(c(ki_kare=unname(ki_kare$statistic),
  df=unname(ki_kare$parameter), p=ki_kare$p.value,
  cramer_v=sqrt(unname(ki_kare$statistic)/sum(tablo)),
  min_beklenen=min(ki_kare$expected),
  bes_alti_hucre=sum(ki_kare$expected < 5)), digits=15)
\end{lstlisting}

\textbf{Çıktıyı okuma.} Paylaşılan R ekranında
$\chi^2=1{,}969808570687404$, $p=0{,}160468182571215$ ve
$V=0{,}070617682919937$ vardır. Yüzdelerin paydaları satır
toplamlarıdır: F için 208, M için 187.
\end{minipage}\par

\par\noindent\begin{minipage}{\linewidth}
\textbf{Korelasyon, ayrı alt grup ve grafik.} Ana veri korunur.

\begin{lstlisting}[style=prismR]
korelasyon <- cor.test(veri$g1, veri$g3,
  method="pearson", alternative="two.sided")
print(c(n=nrow(veri), r=unname(korelasyon$estimate),
  p_iki_yonlu=korelasyon$p.value), digits=15)
pozitif <- veri[veri$g3 > 0, ]
print(rbind(Pozitif_G3=betimsel(pozitif$g3),
            Tum_G3=betimsel(veri$g3)), digits=15)
stopifnot(nrow(veri) == 395, nrow(pozitif) == 357)
par(mfrow=c(1, 1))
boxplot(veri$g3, main="Tum kayitlarda yil sonu notu",
        ylab="G3")
\end{lstlisting}

\textbf{Çıktıyı okuma.} R'de $r=0{,}801467932017414$,
$p\approx9{,}00143\times10^{-90}$; pozitif notlarda 357 kayıt ve
11,5238095238095 ortalama görülür. R'nin sayısal ekranları ve kutu
grafiği incelenmiştir; R bu ortamda yeniden çalıştırılmamıştır.
\end{minipage}\par

\par\noindent\begin{minipage}{\linewidth}
\subsection{Python uygulaması}
\label{sec:python-b14}

Blokları sırayla çalıştırın. \texttt{pozitif} ayrı bir kopyadır;
bu seçim ana testlerden önce uygulanmaz. \texttt{ddof=1} örneklem
standart sapmasını, \texttt{.15g} ise 15 anlamlı basamağı seçer.

\begin{lstlisting}[style=prismPythonApp]
import pandas as pd
import numpy as np
from scipy import stats

veri = pd.read_csv("companion/spss/csv/b14.csv")
assert veri.shape == (395, 10)
assert not veri.isna().any().any()
assert veri.id.is_unique
assert veri.sex.isin([1, 2]).all()
assert veri.pass10.isin([0, 1]).all()
assert veri.pass10.eq((veri.g3 >= 10).astype(int)).all()

def betimsel(degerler):
    return dict(n=len(degerler), ortalama=degerler.mean(),
        ss=degerler.std(ddof=1), minimum=degerler.min(),
        maksimum=degerler.max())

def yazdir(baslik, ozet):
    print("\n" + baslik)
    for ad, deger in ozet.items():
        print(f"{ad}: {deger:.15g}")

def tablo_yazdir(baslik, cerceve):
    print("\n" + baslik)
    print(cerceve.to_string(
        float_format=lambda deger: f"{deger:.15g}"))
\end{lstlisting}
\end{minipage}\par

\par\noindent\begin{minipage}{\linewidth}
\textbf{Veri kontrolü ve ortalama testi.} Aralıkların yönü aynıdır.

\begin{lstlisting}[style=prismPythonApp]
yazdir("Veri kontrolu", dict(satir=len(veri),
    sutun=veri.shape[1], eksik=int(veri.isna().sum().sum()),
    sifir_not=int((veri.g3 == 0).sum())))
yazdir("G1 betimsel", betimsel(veri.g1))
yazdir("G3 betimsel", betimsel(veri.g3))
test = stats.ttest_1samp(veri.g3, 10,
                         alternative="two-sided")
aralik = test.confidence_interval(.95)
yazdir("Tek orneklem testi ve guven araliklari", dict(
    t=test.statistic, df=len(veri)-1,
    p_iki_yonlu=test.pvalue,
    ortalama_fark=veri.g3.mean()-10,
    ortalama_alt=aralik.low, ortalama_ust=aralik.high,
    fark_alt=aralik.low-10, fark_ust=aralik.high-10))
\end{lstlisting}

\textbf{Çıktıyı okuma.} Python çıktısında 395 satır, 10 sütun,
sıfır eksik değer ve 38 sıfır not vardır. Tek örneklem testinin
iki yönlü $p$ değeri R ile aynı yuvarlanmış sonucu verir.
\end{minipage}\par

\par\noindent\begin{minipage}{\linewidth}
\textbf{Çapraz tablo ve Pearson testi.} Düzeltme kapalıdır.

\begin{lstlisting}[style=prismPythonApp]
tablo = pd.crosstab(veri.sex, veri.pass10).reindex(
    index=[1, 2], columns=[0, 1], fill_value=0)
tablo.index = ["F", "M"]
tablo.columns = ["10_alti", "10_ve_uzeri"]
tablo_yazdir("Capraz tablo", tablo)
print("Satir toplamlari:", tablo.sum(axis=1).to_dict())
print("Sutun toplamlari:", tablo.sum(axis=0).to_dict())
print("Genel toplam:", int(tablo.to_numpy().sum()))
yuzdeler = 100*tablo.div(tablo.sum(axis=1), axis=0)
tablo_yazdir("Satir yuzdeleri", yuzdeler)
ki_kare, p_degeri, serbestlik, beklenen = (
    stats.chi2_contingency(tablo, correction=False))
tablo_yazdir("Beklenen sayimlar", pd.DataFrame(
    beklenen, index=tablo.index, columns=tablo.columns))
yazdir("Pearson ki-kare testi", dict(
    ki_kare=ki_kare, df=serbestlik, p=p_degeri,
    cramer_v=np.sqrt(ki_kare/tablo.to_numpy().sum()),
    min_beklenen=beklenen.min(),
    bes_alti_hucre=int((beklenen < 5).sum())))
\end{lstlisting}
\end{minipage}\par

\par\noindent\begin{minipage}{\linewidth}
\textbf{Korelasyon, ayrı alt grup ve grafik.} Tüm kayıtlar korunur.

\begin{lstlisting}[style=prismPythonApp]
korelasyon = stats.pearsonr(veri.g1, veri.g3)
yazdir("Pearson korelasyonu", dict(n=len(veri),
    r=korelasyon.statistic, p_iki_yonlu=korelasyon.pvalue))
pozitif = veri.loc[veri.g3 > 0].copy()
yazdir("Pozitif G3", betimsel(pozitif.g3))
yazdir("Tum G3", betimsel(veri.g3))
assert len(veri) == 395
assert len(pozitif) == 357

import matplotlib.pyplot as plt

plt.figure(figsize=(6, 5))
plt.boxplot(veri.g3)
plt.title("Tum kayitlarda yil sonu notu")
plt.ylabel("G3")
plt.xticks([1], ["Tum kayitlar"])
plt.tight_layout(); plt.show()
\end{lstlisting}

\textbf{Çıktıyı okuma.} Paylaşılan Python ekranlarında pozitif G3
için $n=357$, ortalama 11,5238095238095; tüm G3 için $n=395$,
ortalama 10,4151898734177 görülür. Pozitif alt grupta yeni bir
anlamlılık testi yapılmamıştır; bu adım yalnız betimsel karşılaştırmadır.
\end{minipage}\par

\Needspace{12\baselineskip}
\subsection{Sonuçların karşılaştırılması ve ortak rapor}
13 Eylül 2026'da paylaşılan SPSS tabloları, R ekranları ve tamamlanan
Python terminal çıktıları karşılaştırılmıştır. Betimsel özetler,
tek örneklem testinin sonuçları ve aralıkları, çapraz tablonun sayımları,
yüzdeleri ve beklenen sayımları, Pearson ki-kare ve korelasyon
sonuçları ile pozitif alt grup özetleri uyumludur. Karşılaştırılan
R ve Python değerlerinde $10^{-12}$ mutlak tolerans; yaklaşık
$9\times10^{-90}$ büyüklüğündeki korelasyon $p$ değerinde ise
$10^{-10}$ bağıl tolerans kullanılmıştır. Bütün basamaklarda
eşitlik iddia edilmez. Python kodu proje CSV'siyle ayrıca
çalıştırılmıştır; SPSS değerleri gösterilen yuvarlama hassasiyetlerinde
uyumludur. Üç yazılımın G3 kutu grafikleri de incelenmiştir.
Aşağıdaki tablolar paylaşılan çıktıların kitap düzenindeki özetidir;
ekran görüntüsü değildir.

\begin{table}[H]
\centering
\caption{B14'te tam veri ve pozitif notlu alt grubun özetleri}
\label{tab:b14-dogrulanmis-betimsel}
\begin{tabular}{@{}lrrrr@{}}
\toprule
Değişken/kapsam & $n$ & Ortalama & Standart sapma & Min--Maks \\
\midrule
G1 (tümü) & 395 & 10,908861 & 3,319195 & 3--19 \\
G3 (tümü) & 395 & 10,415190 & 4,581443 & 0--20 \\
G3 (pozitif) & 357 & 11,523810 & 3,227797 & 4--20 \\
\bottomrule
\end{tabular}
\par\smallskip
{\footnotesize Satırlar bağımsız örneklemler olarak toplanmaz.
Pozitif G3, ana verinin alt grubudur. SPSS'in son tablosu tekrar
395 kayıt ve 10,415190 ortalama göstererek geçici seçimin sona
erdiğini doğrulamıştır; R ve Python ana veriyi değiştirmemiştir.\par}
\end{table}

\begin{table}[H]
\centering
\caption{B14'te yıl sonu ortalaması ve 10 puana karşı test}
\label{tab:b14-dogrulanmis-ortalama}
\begin{tabular}{@{}lr@{}}
\toprule
Ölçü & Değer \\
\midrule
Ortalama & 10,415190 \\
Ortalamanın standart hatası & 0,230517 \\
Ortalama için \%95 güven aralığı & $[9{,}961992;\ 10{,}868388]$ \\
10'dan ortalama fark & 0,415190 \\
Fark için \%95 güven aralığı & $[-0{,}038008;\ 0{,}868388]$ \\
$t$ istatistiği & 1,801122 \\
Serbestlik derecesi & 394 \\
İki yönlü $p$ değeri & 0,072448 \\
\bottomrule
\end{tabular}
\par\smallskip
{\footnotesize Her iki aralık aynı 395 kayda dayanır. Fark aralığı,
ortalama aralığının iki sınırından 10 çıkarılarak elde edilir.
Tek yönlü $p$ bu uygulamanın kararında kullanılmaz.\par}
\end{table}

\begin{table}[H]
\centering
\caption{B14'te 10 puan eşiği için sayımlar ve satır yüzdeleri}
\label{tab:b14-dogrulanmis-sayim}
\begin{tabular}{@{}lrrrr@{}}
\toprule
Grup & $<10$ & $\geq10$ & Toplam & $\geq10$ yüzdesi \\
\midrule
F & 75 & 133 & 208 & 63,942308 \\
M & 55 & 132 & 187 & 70,588235 \\
\midrule
Toplam & 130 & 265 & 395 & 67,088608 \\
\bottomrule
\end{tabular}
\par\smallskip
{\footnotesize Tam 10 puan eşiği karşılayan kategoriye dahildir.
Bağımsızlık altında beklenen sayımlar F satırında
$(68{,}455696;\ 139{,}544304)$, M satırında
$(61{,}544304;\ 125{,}455696)$'dır.\par}
\end{table}

\begin{table}[H]
\centering
\caption{B14'te iki ayrı ilişki sorusunun sonuçları}
\label{tab:b14-dogrulanmis-iliski}
\begin{tabular}{@{}lr@{}}
\toprule
Ölçü & Değer \\
\midrule
\texttt{sex}--\texttt{pass10}: Pearson $\chi^2(1)$ & 1,969809 \\
Pearson ki-kare $p$ değeri & 0,160468 \\
Cramér $V$ & 0,070618 \\
En küçük beklenen sayım & 61,544304 \\
Beklenen sayımı 5'in altındaki hücre & 0 \\
\midrule
G1--G3: Pearson $r$ & 0,801468 \\
Korelasyonun iki yönlü $p$ değeri (yaklaşık) & $9{,}00143\times10^{-90}$ \\
\bottomrule
\end{tabular}
\par\smallskip
{\footnotesize İki analiz de 395 kayıt kullanır fakat farklı
hipotezleri sınar. Ki-kare testinde süreklilik düzeltmesi yoktur.
Çoğu değer altı ondalık basamakla gösterilmiştir; ara hesaplarda
yuvarlanmamış değerler kullanılmıştır.\par}
\end{table}

\textbf{Grafik kontrolünün kapsamı.} SPSS, R ve Python'un tüm G3
kayıtlarını içeren kutu grafiklerinde medyan 11, kutu sınırları
8 ve 14, bıyıklar 0 ve 20'dir. Bu toplu grafikte sıfırlar ayrıca
aykırı işaretlenmemiştir. Bu durum B11'de F ve M ayrı çizildiğinde
sıfırların işaretlenmesiyle çelişmez: grup çeyrekleri ve kutu
genişlikleri farklıdır. Aykırı etiketi verinin hatalı olduğunu
göstermez; etiketin olmaması da normalliği veya bağımsızlığı kanıtlamaz.
Kutudaki çizgi ortalamayı değil medyanı gösterir. B14'te incelenen
grafikler kutu grafikleridir; bunlardan korelasyonun doğrusal
örüntüsünün kontrol edildiği sonucu çıkarılmaz.

\Needspace{12\baselineskip}
\begin{outputguidebox}
Tablo~\ref{tab:b14-dogrulanmis-ortalama} için $p=0{,}072448>0{,}05$
olduğundan 10 puana eşit ortalama sıfır hipotezi reddedilmez;
bu, ortalamanın kesinlikle 10 olduğu anlamına gelmez.
Tablo~\ref{tab:b14-dogrulanmis-iliski} içindeki kategorik analizde
de bağımsızlık hipotezi reddedilmez; bağımsızlık kanıtlanmış olmaz.
Pozitif G1--G3 korelasyonu ise farklı bir sorunun bulgusudur.
Anlamlı sonuçları seçip diğerlerini gizlemek yerine her sonucu
kendi araştırma sorusuyla birlikte raporlayın.

Tablo~\ref{tab:b14-dogrulanmis-betimsel} içindeki pozitif G3
ortalamasının artması, sonuca göre seçimin etkisini gösterir.
Bu dosyada 395 kaydın 38'inin notu sıfır olduğundan, tüm verinin
ortalaması pozitif alt grubun ortalamasının $357/395$ katıdır.
Bu aritmetik ilişki, dışlamanın doğruluğuna ilişkin kanıt değildir.
Buradaki duyarlılık özeti aynı hedef için alternatif bir test
değil, farklı bir alt grubun betimlemesidir; 357 ve 395 kayıt
bağımsız iki grup gibi test edilmez. Sıfırların nasıl oluştuğu
kaynak bilgisiyle değerlendirilmeden eksik değer oldukları varsayılmaz.

Yazılımların uyumu gözlemlerin bağımsızlığını veya iki okulla
sınırlı verinin temsil gücünü doğrulamaz. Testler okul ve diğer
karıştırıcılar için düzeltme içermez. Korelasyon nedensel etki
değildir; 0,05 düzeyindeki kararlar bilimsel önemle eş tutulmaz.
B10, B12 ve B13 ile aynı kayıtların yeniden analizi bağımsız
bir tekrar çalışma sayılmaz.
\end{outputguidebox}

\Needspace{10\baselineskip}
\textbf{Raporlama örneği (doğrulanmış çıktılarla).}
``Ana analizlerde 395 öğrencinin tüm kayıtları kullanılmış ve
sıfır yıl sonu notu bulunan 38 kayıt korunmuştur. G3 ortalaması
10,415190, standart sapması 4,581443'tür (ortalama için \%95 güven
aralığı $[9{,}961992;\ 10{,}868388]$). 10 puana karşı iki yönlü
testte sıfır hipotezi reddedilmemiştir
($t(394)=1{,}801122$, $p=0{,}072448$). Kaynak cinsiyet kategorisi
ile 10 puan eşiğini karşılama durumu için süreklilik düzeltmesiz
Pearson testinde yeterli ilişki kanıtı elde edilmemiştir
($\chi^2(1)=1{,}969809$, $p=0{,}160468$, Cramér $V=0{,}070618$).
Ayrı doğrusal ilişki sorusunda G1--G3 korelasyonu 0,801468'dir
($p\approx9{,}00143\times10^{-90}$). Her test ayrı 0,05 düzeyinde
değerlendirilmiş; çoklu test düzeltmesi yapılmamıştır.
Yalnız pozitif notları içeren betimsel duyarlılık özetinde
357 kaydın ortalaması 11,523810 bulunmuştur; bu, ana analizlerin
yerine geçmeyen farklı bir alt grup özetidir. Grafik ve yazılım
uyumu varsayımları doğrulamamış; okul ve diğer karıştırıcılar
denetlenmemiştir. Bulgular nedensel etki veya daha geniş bir evrene
genellenebilirlik kanıtı olarak sunulmamıştır.''

\textbf{Raporlama görevi.} Veri kaynağı ve lisansı, gözlem birimi,
seçilen değişkenler, betimsel özet, soruya uygun tek temel analiz
ve sınırlılıkları içeren bir sayfalık rapor yazın. Temel analizi
$p$ değerine göre seçmeyin; diğer soruları ayrı ve açık biçimde
belirtin. Sıfırların nasıl ele alındığını ve pozitif alt grubun
neden farklı bir hedef olduğunu açıklayın. Ortalama aralığını
fark aralığıyla, tek yönlü $p$ değerini iki yönlü değerle
karıştırmayın. Küçük $p$ değerlerini sıfır olarak yazmayın.
Çalıştırılan kodları ve yazılımların ürettiği çıktıları raporla saklayın.

\section{R ile çözümlü uygulama}
\label{sec:r-b14}

\textbf{Soru.} Aynı 24 kişinin son--ön farkları için $\bar d=-3{,}2$ ve
$s_d=6$ verilmiştir. Uygun testi seçin, yüzde 95 güven aralığını ve
standartlaştırılmış değişimi hesaplayıp sonucu raporlayın.

\begin{lstlisting}[style=prismR]
cift_sayisi <- 24
ortalama_fark <- -3.2
fark_sapmasi <- 6
standart_hata <- fark_sapmasi / sqrt(cift_sayisi)
serbestlik <- cift_sayisi - 1
t_degeri <- ortalama_fark / standart_hata
p_degeri <- 2 * pt(abs(t_degeri), df = serbestlik,
                   lower.tail = FALSE)
aralik <- ortalama_fark + c(-1, 1) *
  qt(0.975, df = serbestlik) * standart_hata
sonuc <- data.frame(
  fark = ortalama_fark, se = standart_hata,
  t = t_degeri, sd = serbestlik, p = p_degeri,
  alt = aralik[1], ust = aralik[2],
  dz = ortalama_fark / fark_sapmasi
)
print(sonuc, digits = 5)
stopifnot(p_degeri < 0.05, all(aralik < 0))
\end{lstlisting}

\begin{outputguidebox}
\textbf{Çözüm ve kontrol değerleri.} Tasarım eşleştirilmiştir; farklar
üzerinden tek örneklem $t$ hesabı yapılır. $\SE=1{,}2247$,
$t(23)=-2{,}6128$, $p=0{,}01556$, yüzde 95 güven aralığı
$[-5{,}7336;-0{,}6664]$ ve $d_z=-0{,}5333$ bulunur.

\textbf{Örnek rapor.} ``Son ölçümde puanlar ortalama 3,2 puan daha düşüktü;
$t(23)=-2{,}61$, $p=0{,}016$, yüzde 95 GA $[-5{,}73;-0{,}67]$,
$d_z=-0{,}53$.''
\end{outputguidebox}

\textbf{Yorum.} 24 bağımsız kişi çifti vardır; 48 bağımsız kişi yoktur.
Fark dağılımındaki aykırı değerler ve yaklaşık normallik ham veriden
incelenmelidir. Azalmanın olumlu mu olumsuz mu olduğu ölçeğin anlamına bağlıdır;
kontrol grubu yoksa değişim yalnız uygulamaya bağlanamaz.

\textbf{Pekiştirme ve yanıt.} Son ölçümden ön ölçüme geçildiğinde farkın
ve $t$'nin işareti değişir, çift yönlü $p$ aynı kalır. Sınav yanıtında yöntem,
işaret, belirsizlik ve bağlam birlikte verilmelidir.

\begin{etkinlik}
Rastgele bir soruyu seçip yalnız sonucunu değil karar izini yazın: hedef
büyüklük, veri yapısı, yöntem, varsayım, tahmin, belirsizlik, etki büyüklüğü
ve ileri sürülemeyecek iddia.
\end{etkinlik}

\section{Bölüm özeti}

\begin{itemize}
  \item Yöntem seçimi araştırma hedefi, değişken türü ve gözlemlerin ilişkisine dayanır.
  \item Tam bir çözüm hedef parametreyi, varsayımları, hesabı ve bağlamsal yorumu içerir.
  \item \pval-değeri tek başına etki büyüklüğü, hassasiyet veya nedensellik göstermez.
  \item Yazılım hesapları doğrular; bilimsel kararın gerekçesini araştırmacı kurar.
  \item Aynı sayısal çıktı, örnekleme ve araştırma tasarımı farklı olduğunda
  farklı genelleme ve nedensellik sınırlarına sahiptir.
  \item Güven aralığı ve etki büyüklüğü, ikili anlamlı/anlamsız kararından
  daha zengin bir sonuç değerlendirmesi sağlar.
\end{itemize}

\bolumkaynaklari{\cite{gaise2016,cumming2014,appelbaum2018}}

\section*{Alıştırmalar}

\subsection*{Genel tekrar soruları}

\begin{enumerate}
  \item Örnekleme yanlılığı büyük örneklemle neden kendiliğinden kaybolmaz?
  \item MLT'nin istatistiksel çıkarımdaki rolünü açıklayın.
  \item Yüzde 95 güven aralığını doğru biçimde yorumlayın.
  \item Tip I hata, Tip II hata ve test gücü arasındaki ilişkiyi açıklayın.
  \item Bağımsız ve eşleştirilmiş \tstat testlerini tasarım bakımından karşılaştırın.
  \item Ki-kare bağımsızlık testinin koşullarını yazın.
  \item Güçlü korelasyonun neden nedensellik kanıtı olmadığını açıklayın.
  \item Bir regresyon eğimi ve $R^2$ değeri için raporlama cümlesi kurun.
  \item Standart sapma ile standart hatayı tek bir araştırma örneği üzerinde karşılaştırın.
  \item İstatistiksel anlamlılık ile uygulamalı önemin ayrıldığı bir sonuç tasarlayın.
  \item Welch ve eşleştirilmiş \tstat testlerinin analiz birimlerini karşılaştırın.
  \item Ki-kare artıklarının ve Cramér $V$ değerinin farklı rollerini açıklayın.
\end{enumerate}

\subsection*{Mini deneme}

\begin{enumerate}
  \item Bir üniversitedeki tüm birinci sınıf öğrencileri araştırmanın hedef
  grubudur. Bu grubun yalnız gönüllü 80 üyesi incelenmiştir. Evreni, örneklemi
  ve en önemli yanlılık riskini belirtin.
  \item Standart sapması 15 olan bir değişkende $n=100$ için ortalamanın
  standart hatasını hesaplayın.
  \item $\bar x=68$, $s=12$, $n=36$ için yaklaşık yüzde 95 güven aralığını
  kritik değeri 2 alarak hesaplayın.
  \item Bir testte $p=0{,}08$ ve $\alpha=0{,}05$ bulunmuştur. Uygun kararı ve
  kullanılması gereken ifadeyi yazın.
  \item Aynı öğrencilerin dönem başı ve dönem sonu puanları hangi \tstat testiyle
  karşılaştırılır? Analizde hangi puanların dağılımı incelenmelidir?
  \item Bir çapraz tabloda tüm beklenen frekanslar 5'in altında kalmıştır.
  Pearson ki-kare yaklaşımı bakımından ne yapılmalıdır?
  \item $r=-0{,}60$ değerinin yönünü, gücünü ve basit doğrusal modeldeki
  $R^2$ karşılığını yorumlayın.
  \item Anlamlı fakat çok dar uygulamalı etkiye sahip bir bulgunun neden
  yalnız \pval-değeriyle raporlanmaması gerektiğini açıklayın.
  \item $n=64$, $\bar x=75$, $s=16$ için standart hatayı ve kritik değeri 2
  alarak yaklaşık yüzde 95 güven aralığını bulun.
  \item $n=36$, $\bar x=52$, $s=12$ ve $\mu_0=48$ için tek örneklem t
  değerini ve Cohen'in $d$ değerini hesaplayın.
  \item İki bağımsız grupta ortalama farkı 5 ve farkın standart hatası 2 ise
  t değerini ve yaklaşık yüzde 95 güven aralığını bulun.
  \item $2\times3$ çapraz tabloda serbestlik derecesini hesaplayın ve anlamlı
  genel testten sonra hangi iki çıktının incelenmesi gerektiğini yazın.
  \item $n=40$ ve $r=0{,}45$ için $R^2$ değerini hesaplayıp yorumlayın.
  \item Bir artık grafiğinde belirgin huni biçimi görülmesinin hangi varsayımı
  sorgulattığını açıklayın.
\end{enumerate}

\subsection*{Mini deneme yanıtları}

\begin{enumerate}
  \item Evren tüm birinci sınıf öğrencileri, örneklem 80 gönüllüdür; gönüllü
  katılım nedeniyle öz-seçim yanlılığı temel risktir.
  \item $SE=15/\sqrt{100}=1{,}5$.
  \item $68\pm2(12/6)=68\pm4$, yani yaklaşık $[64;72]$.
  \item $H_0$ reddedilemez; ``fark yoktur'' yerine ``$H_0$'ı reddetmek için
  yeterli kanıt bulunamamıştır'' denir.
  \item Eşleştirilmiş örneklem \tstat testi kullanılır; son--ön fark puanlarının
  dağılımı değerlendirilir.
  \item Kesin yöntemler, kategori birleştirmenin bilimsel olarak anlamlı olup
  olmadığı veya daha büyük örneklem seçenekleri değerlendirilmelidir.
  \item Orta-güçlü negatif doğrusal ilişki vardır; $R^2=0{,}36$'dır.
  \item Etki tahmini ve güven aralığı, bulgunun büyüklüğünü ve hassasiyetini
  gösterir; küçük \pval-değeri tek başına uygulamalı önem göstermez.
\end{enumerate}